%====================================================
% CHAPTER I
%====================================================

\chapter{Project Introduction}

\section{Overview}

\subsection{Project Information}

\begin{itemize}
    \item \textbf{Project name:} Personalized Career Orientation \& Learning Roadmap Platform for Software Engineering Students
    \item \textbf{Project code:}
    \item \textbf{Group name:} Group 2
    \item \textbf{Software Type:} Web application
\end{itemize}

The rapid development of Information Technology has significantly transformed the way students learn, develop professional skills, and prepare for future careers. In particular, Software Engineering students are required to continuously update their technical knowledge, practical experience, and industry trends to remain competitive in the job market.

However, many students face difficulties in determining a suitable career path, identifying the required technical competencies, evaluating their current skill levels, and planning an effective learning roadmap. They often rely on scattered learning resources, inconsistent career advice, or personal experiences without systematic guidance. Consequently, students may spend considerable time learning technologies that are not aligned with their career goals or current industry demands.

The \textbf{Personalized Career Orientation \& Learning Roadmap Platform for Software Engineering Students} is proposed to address these challenges. The platform provides intelligent career guidance by analyzing students’ academic background, technical skills, career interests, and learning progress. Based on this information, the system generates personalized learning roadmaps, recommends appropriate learning resources, identifies missing competencies through skill gap analysis, and helps students build professional portfolios using their GitHub projects.

By integrating Artificial Intelligence, career planning, learning management, portfolio generation, and job market analysis into a unified platform, the proposed system aims to improve students' learning efficiency and enhance their career readiness.


\subsection{Project Team}

\begin{table}[H]
\centering
\caption{Project Team}
\begin{tabular}{|c|c|c|c|}
\hline
\textbf{Full Name} &
\textbf{Role} &
\textbf{Email} &
\textbf{Mobile} \\
\hline

Nguyen Van Bao &
Project Manager &
&
\\
\hline

Phan Thi Duy Chung &
Business Analyst &
&
\\
\hline

Vo Thi Kiem Dung &
System Analyst &
&
\\
\hline

Tran Thi Bich Diem &
UI/UX Designer &
&
\\
\hline

Vo Trong Nguyen &
Backend Developer &
&
\\
\hline

Vu Le Hoang Nhat &
Frontend Developer &
&
\\
\hline

Vo Bao Phuc &
Tester &
&
\\
\hline

\end{tabular}
\end{table}


\section{Product Background}

Software Engineering education focuses not only on theoretical knowledge but also on practical skills, project experience, and continuous learning. Modern software industries require graduates to possess diverse competencies, including programming, database management, software architecture, cloud computing, DevOps, artificial intelligence, and communication skills.

Despite the availability of numerous online learning platforms, students often encounter several challenges:

\begin{itemize}
    \item Lack of personalized career guidance.
    \item Difficulty selecting an appropriate Software Engineering specialization.
    \item Limited understanding of industry-required technical skills.
    \item Inability to evaluate their own competencies objectively.
    \item Difficulty organizing learning plans according to career objectives.
    \item Lack of professional portfolios for internship and job applications.
\end{itemize}

Currently, students usually collect learning materials from various websites such as Microsoft Learn, Coursera, Udemy, roadmap.sh, GitHub, YouTube, and Stack Overflow. Since these resources are independent and not integrated, students must manually determine what to study, when to study, and how to evaluate their progress.

Therefore, a comprehensive platform that combines career orientation, personalized learning roadmap generation, AI-assisted mentoring, skill assessment, portfolio management, and labor market analysis can significantly improve students' learning outcomes and employment opportunities.


%====================================================
\section{Existing Systems}

Several existing platforms provide partial support for career planning or online learning; however, none of them fully integrate all required functionalities for Software Engineering students. 

\textbf{Roadmap.sh} 
    
    Roadmap.sh provides structured technology roadmaps for different software engineering career paths, including Backend Development, Frontend Development, DevOps, AI Engineering, and Mobile Development. Although it offers detailed technology trees, it does not personalize learning paths based on students' individual competencies or academic performance. 
    
\textbf{Coursera} 

    Coursera offers thousands of online courses from universities and technology companies. Students can obtain certificates after completing courses. However, course recommendations are primarily based on browsing history rather than comprehensive career analysis. 
    
\textbf{LinkedIn Learning} 

    LinkedIn Learning provides professional courses and integrates with LinkedIn profiles. Nevertheless, it lacks personalized skill gap analysis and automatic roadmap generation tailored to Software Engineering students. 

\textbf{GitHub} 

    GitHub serves as the largest source code hosting platform, allowing students to manage software projects and demonstrate programming experience. However, GitHub itself does not evaluate students' technical competencies or generate learning recommendations. 

\textbf{ChatGPT and Gemini} 

    Artificial Intelligence chatbots such as ChatGPT and Gemini can answer programming and career-related questions. Nevertheless, they do not maintain long-term learning progress or automatically generate structured career roadmaps integrated with students' academic profiles. 

    Therefore, there is still a need for an integrated platform specifically designed for Software Engineering students. 

\section{Business Opportunity}
The software industry continues to experience rapid growth worldwide. Companies increasingly demand graduates with practical technical skills rather than solely academic qualifications. 

Universities are also shifting toward outcome-based education, emphasizing competency assessment, project-based learning, and career preparation. Consequently, educational institutions require systems capable of supporting students throughout their learning journeys. 

The proposed platform provides several business opportunities: 

\begin{itemize}
    \item Supporting universities in career counseling.  
    \item Assisting students in planning personalized learning paths.  
    \item Helping employers evaluate candidates through dynamic portfolios.  
    \item Assisting academic advisors in monitoring students' learning progress.  
    \item Providing technology trend analysis based on labor market demands.  
    \item Integrating Artificial Intelligence into educational support systems.  
\end{itemize}

Furthermore, the platform can be expanded in the future to support multiple universities, online certification providers, internship management systems, and enterprise recruitment platforms. 

\section{Software Product Vision}
The vision of the project is to become an intelligent career development platform that accompanies Software Engineering students throughout their academic journey and professional development. 

The system aims to: 

\begin{itemize}
    \item Provide personalized career recommendations.  
    \item Generate intelligent learning roadmaps.  
    \item Continuously analyze students' skill gaps.
    \item Recommend high-quality learning resources.  
    \item Build professional online portfolios automatically.  
    \item Analyze current technology trends and employment demands.
    \item Improve students' employability and career readiness through data-driven recommendations.  
\end{itemize}

By integrating Artificial Intelligence, GitHub repositories, academic information, and industry trends, the platform creates a comprehensive ecosystem for career development. 

\section{Project Scope \& Limitations}

\subsection{Major Features}

\subsubsection{Major Features for Administrators}

FE-01: Login/Logout. 

FE-02: Manage user accounts (Students, Academic Counselors, Industry Mentors). 

FE-03: Manage career paths (Backend Developer, Frontend Developer, Full Stack Developer, DevOps Engineer, AI Engineer, Data Engineer, Mobile Developer, QA Engineer, etc.). 

FE-04: Manage skill categories and technical skills. 

FE-05: Manage learning resources (courses, books, documentation, videos, practice websites). 

FE-06: Manage AI prompt templates and system configurations. 

FE-07: Manage job market data and technology trend information. 

FE-08: Manage notifications and announcements. 

FE-09: View system reports and analytics. 

FE-10: Monitor user activities and system logs. 

FE-11: Manage role-based access permissions. 

\subsubsection{Major Features for Student}

FE-01: Register, login, logout, and authenticate using Google OAuth. 

FE-02: Manage personal profile, including academic information, technical skills, career interests, and learning goals. 

FE-03: Complete skill assessment to evaluate current competencies. 

FE-04: Receive AI-powered career consultation through the AI Virtual Mentor. 

FE-05: Select a target software engineering career path. 

FE-06: Generate a personalized learning roadmap based on selected career goals. 

FE-07: View and update learning progress for each roadmap milestone. 

FE-08: Perform skill gap analysis to identify missing competencies. 

FE-09: Receive recommended learning resources based on skill gaps. 

FE-10: Connect and synchronize GitHub repositories. 

FE-11: Automatically generate and maintain an online professional portfolio. 

FE-12: View job market trends, technology demand, and career insights. 

FE-13: Receive notifications about roadmap updates, learning recommendations, and career opportunities. 

\subsubsection{Major Features for Academic Counselor}

FE-01: Login/Logout. 

FE-02: View student profiles and academic information. 

FE-03: Monitor students’ learning roadmap progress. 

FE-04: Review students’ skill gap analysis reports. 

FE-05: Provide academic advice and learning recommendations. 

FE-06: Track students’ learning performance and completion status. 

FE-07: Send notifications or reminders to students. 

FE-08: View reports summarizing students’ learning progress. 

\subsubsection{Major Features for Industry Mentor}

FE-01: Login/Logout. 

FE-02: Review students’ career roadmaps. 

FE-03: Evaluate students’ professional portfolios. 

FE-04: Recommend industry-required technical skills. 

FE-05: Provide career advice and industry feedback. 

FE-06: Suggest learning priorities according to current market demand. 

FE-07: Review students’ project experiences and GitHub repositories. 

FE-08: View students’ competency development reports. 

\subsection{Limitations \& Exclusions}
LI-01: The first version of the system is designed exclusively for Software Engineering students and does not currently support students from other academic disciplines. 

LI-02: AI career recommendations depend on external AI services (such as OpenAI or Gemini). Recommendation quality may vary depending on the AI model and available information. 

LI-03: GitHub integration supports only repositories that are publicly accessible or explicitly authorized by users. 

LI-04: Job market analysis depends on external recruitment platforms and publicly available datasets. Real-time accuracy cannot be fully guaranteed. 

LI-05: The system currently supports predefined Software Engineering career paths only. Users cannot create custom career paths. 

LI-06: The platform is implemented as a web application. Native Android and iOS applications are not included in this version. 

LI-07: The platform requires a stable Internet connection because several core services depend on external APIs (AI services, Google Authentication, GitHub API, and job market data). 

LI-08: The current version does not provide online coding practice, automatic code assessment, or integrated programming environments.

